\clearpage
\section{Especificação formal M12.1: posts, comentários e moderação}
\label{sec:formal-m12-1}
\sloppy

M12.1 formaliza o alvo normativo de Post, Comentário, edição, moderação/remoção
lógica, históricos e leitura com máscara (RF19, RF20, UI-11, Q08/Q10/Q11/Q13),
na composição com M7 (idempotência), M9 (autorização no commit) e M11 (vínculo
canônico/status da Turma). A evidência cobre somente o modelo formal declarado.

\subsection*{CUE}
O contrato CUE (\texttt{\#M12\_1Contrato}) verifica shapes de Post, Comentário,
históricos de edição/moderação, projeção de leitura, efeito mínimo de
notificação e identidade de comando, com enums fechados, nulabilidade,
condicionais (remoção lógica, edição, moderação, tipo de histórico, visão de
leitura) e limites: título 1--150, descrição 1--10\,000, texto 1--2\,000. CUE
não prova concorrência nem ACL de Roteiros (M12.2).

\subsection*{Autoria, participação e turma arquivada}
Post pertence à Turma e ao professor dono; Comentário pertence ao Post e ao
autor. A participação acadêmica é o vínculo canônico atual (M11), nunca o
espelho. Somente o professor dono cria/edita Post; o autor edita o próprio
Comentário; o Chefe intervém apenas por Q13 (moderação/remoção com motivo e
auditoria), sem assumir autoria nem criar Post. Turma \texttt{Arquivada} nega
toda escrita acadêmica, inclusive moderação e para o Chefe (Q08).

\subsection*{Edição, remoção lógica e históricos}
Editar preserva a autoria e grava evento imutável em
\texttt{Historico\_Posts\_Turma}/\texttt{Historico\_Comentario}
(\texttt{edicao}/\texttt{moderacao}) com operador e timestamp. Remover Post é
lógico: preserva documento e histórico (RF25) e o oculta dos colegas, mantendo-o
acessível ao professor dono e à chefia. Nenhuma transição remove histórico
anterior; não há exclusão física nem estado de restauração.

\subsection*{Moderação e máscara de leitura}
Moderar Comentário grava \texttt{moderado}, motivo, moderador e instante sem
apagar o original. A lista autorizada revalida vínculo/estado: o autor vê o
original marcado, colegas veem aviso institucional e professor dono/Chefe veem
original e histórico. A edição posterior pelo autor não desfaz a moderação.
\texttt{ColegaNaoVeOriginalModerado}, \texttt{AutorVeOriginalMarcado} e
\texttt{AuditorVeOriginal} fixam a máscara; leitura exige vínculo, autoria,
ownership ou chefia.

\subsection*{Fronteira M12.2, notificação e composição}
O snapshot \texttt{roteiro\_anexo} é fronteira: publicar com roteiro exige o
predicado abstrato \texttt{acessoRoteiroValidado}; a ACL/Storage/URL de Roteiros
não são provadas aqui. A \texttt{POST}/\texttt{COMENTARIO} entra apenas como
efeito mínimo de notificação, sem conteúdo protegido e sem autorizar leitura; a
caixa, fan-out, expiração e marcação de lida são M13. A composição M7 impede
retry de duplicar histórico/notificação e reuso incompatível de herdar receipt;
M9 faz a revogação de vínculo impedir o commit posterior.

\subsection*{Responsabilidade das camadas e limites}
CUE valida shapes e limites. Alloy valida autoria, participação canônica,
turma arquivada somente leitura, edição/remoção/histórico, moderação e máscara,
fronteira abstrata de roteiro e composição M7/M9. Rust valida proveniência,
ordem exata do receipt e as regras determinísticas de limite, exclusividade de
edição e máscara. A evidência \textbf{não} certifica \texttt{functions/},
frontend, Firestore/Storage Rules, Auth, Storage, a ACL completa de Roteiros
(M12.2), a caixa de notificações (M13) nem concorrência real.

\begingroup\small\raggedright
% Gerado por lcqui-spec-doc; não editar.
\subsection*{M12\_1\_Posts\_Comentarios --- formalização executável M12.1}
Shapes estruturais de Post, Comentário, históricos de edição/moderação, projeção de leitura com máscara, efeito mínimo de notificação e identidade de comando. CUE verifica campos, enums fechados, nulabilidade, condicionais (remoção lógica, edição, moderação, tipo de histórico, visão de leitura) e limites (título 1..150, descrição 1..10000, texto 1..2000, nome de arquivo 1..150, tamanho positivo). Autoridade de participação (vínculo canônico M11), transições, turma arquivada somente leitura, edição concorrente, idempotência M7, autorização M9 e a fronteira abstrata de acesso ao roteiro são verificadas em Alloy; a canonicalização determinística por Rust. CUE não prova concorrência temporal nem ACL de Roteiros (M12.2).
\begin{description}
\item[id\_turma] Tipo: string. Obrigatório: sim. Aceita nulo: não.
SQL: TEXT.
Turma do Post/Comentário; vínculo canônico atual decide participação (M11).
\item[id\_professor] Tipo: string. Obrigatório: sim. Aceita nulo: não.
SQL: TEXT.
Autor imutável do Post.
\item[id\_post] Tipo: string. Obrigatório: sim. Aceita nulo: não.
SQL: TEXT.
Post pai do Comentário/histórico.
\item[id\_usuario] Tipo: string. Obrigatório: sim. Aceita nulo: não.
SQL: TEXT.
Autor imutável do Comentário.
\item[titulo] Tipo: string. Obrigatório: sim. Aceita nulo: não.
SQL: TEXT.
Título do Post; de 1 a 150 caracteres.
\item[descricao] Tipo: string. Obrigatório: sim. Aceita nulo: não.
SQL: TEXT.
Descrição do Post; de 1 a 10000 caracteres (Q10).
\item[texto] Tipo: string. Obrigatório: sim. Aceita nulo: não.
SQL: TEXT.
Texto do Comentário; de 1 a 2000 caracteres (Q10), escapado na renderização.
\item[editado] Tipo: bool. Obrigatório: sim. Aceita nulo: não.
SQL: BOOLEAN.
Marcado na edição; a versão anterior fica no histórico.
\item[editado\_em] Tipo: string. Obrigatório: sim. Aceita nulo: sim.
SQL: TIMESTAMP.
Instante do servidor da última edição; nulo quando não editado.
\item[removido\_da\_apresentacao] Tipo: bool. Obrigatório: sim. Aceita nulo: não.
SQL: BOOLEAN.
Remoção lógica de Post (RF25); não apaga documento nem histórico.
\item[motivo\_remocao] Tipo: string. Obrigatório: sim. Aceita nulo: sim.
SQL: TEXT.
Obrigatório quando removido\_da\_apresentacao; até 2000 caracteres.
\item[moderado] Tipo: bool. Obrigatório: sim. Aceita nulo: não.
SQL: BOOLEAN.
Comentário moderado; a moderação não apaga o texto original.
\item[motivo\_moderacao] Tipo: string. Obrigatório: sim. Aceita nulo: sim.
SQL: TEXT.
Obrigatório quando moderado; até 2000 caracteres.
\item[tipo\_historico] Tipo: enum. Obrigatório: sim. Aceita nulo: não.
SQL: ENUM.
Valores: edicao, moderacao.
\item[visao] Tipo: enum. Obrigatório: sim. Aceita nulo: não.
SQL: ENUM.
Valores: AUTOR, COLEGA, AUDITOR.
\item[contem\_conteudo\_protegido] Tipo: bool. Obrigatório: sim. Aceita nulo: não.
SQL: BOOLEAN.
Efeito de notificação nunca transporta conteúdo protegido.
\item[id\_operacao] Tipo: string. Obrigatório: sim. Aceita nulo: não.
SQL: TEXT.
Identidade de comando M7 (uid, tipo\_operacao, payload\_hash).
\item[storage\_path] Tipo: string. Obrigatório: sim. Aceita nulo: não.
SQL: TEXT.
Localizador do snapshot do roteiro; download mediado por endpoint (fronteira M12.2).
\end{description}

% Gerado por lcqui-spec-doc; não editar.
\subsection*{Evidência formal M12.1 --- posts, comentários, edição/moderação e históricos}
CUE verifica shapes de Post, Comentário, históricos de edição/moderação, projeção de leitura com máscara, efeito mínimo de notificação e identidade de comando, com enums fechados, nulabilidade, condicionais (remoção lógica, edição, moderação, tipo de histórico, visão de leitura) e limites (título 1..150, descrição 1..10000, texto 1..2000). Alloy verifica a autoria imutável, a participação por vínculo canônico atual OU professor dono (M11; claim/cache/espelho não recriam participação, de modo que ex-aluno com claim atualizada não edita nem lê), turma arquivada somente leitura (Q08) inclusive para o Chefe por propriedade transversal, criação/edição de Post só pelo dono, criação/edição de Comentário só pelo autor participante, moderação por professor dono ou Chefe Q13, remoção lógica de Post preservando documento/histórico, histórico imutável de edição e moderação, a máscara de leitura (autor vê original marcado, colega vê aviso, auditor vê original), a fronteira abstrata de acesso ao roteiro (sem ACL de Storage), o efeito de notificação vinculado ao idOperacao e sem conteúdo nem autorização de leitura, e a composição M7 (primeira execução grava receipt sob o idOperacao; retry idêntico devolve o resultado e não reaplica; reuso incompatível é rejeitado sem herdar receipt) e M9 (revogação impede commit; participação revalidada no commit). Rust valida proveniência, ordem exata do receipt e as regras determinísticas de limite, exclusividade de edição e máscara. A evidência não certifica Firebase, backend, Rules, Storage, ACL de Roteiros (M12.2), caixa de notificações (M13) nem concorrência real.
\begin{description}
\item[M12\_1-INV-001] CriarPostSoEmAtivo: UNSAT.\newline Escopo: \texttt{check CriarPostSoEmAtivo for 6}.
\item[M12\_1-INV-002] EditarPostSoEmAtivo: UNSAT.\newline Escopo: \texttt{check EditarPostSoEmAtivo for 6}.
\item[M12\_1-INV-003] RemoverPostSoEmAtivo: UNSAT.\newline Escopo: \texttt{check RemoverPostSoEmAtivo for 6}.
\item[M12\_1-INV-004] CriarComentSoEmAtivo: UNSAT.\newline Escopo: \texttt{check CriarComentSoEmAtivo for 6}.
\item[M12\_1-INV-005] EditarComentSoEmAtivo: UNSAT.\newline Escopo: \texttt{check EditarComentSoEmAtivo for 6}.
\item[M12\_1-INV-006] ModerarComentSoEmAtivo: UNSAT.\newline Escopo: \texttt{check ModerarComentSoEmAtivo for 6}.
\item[M12\_1-INV-007] TurmaArquivadaNegaEscrita: UNSAT.\newline Escopo: \texttt{check TurmaArquivadaNegaEscrita for 6}.
\item[M12\_1-INV-008] ChefeNaoCriaPost: UNSAT.\newline Escopo: \texttt{check ChefeNaoCriaPost for 6}.
\item[M12\_1-INV-009] ChefeNaoEditaConteudo: UNSAT.\newline Escopo: \texttt{check ChefeNaoEditaConteudo for 6}.
\item[M12\_1-INV-010] AutorImutavelPost: UNSAT.\newline Escopo: \texttt{check AutorImutavelPost for 6}.
\item[M12\_1-INV-011] AutorImutavelComent: UNSAT.\newline Escopo: \texttt{check AutorImutavelComent for 6}.
\item[M12\_1-INV-012] TerceiroNaoEditaPost: UNSAT.\newline Escopo: \texttt{check TerceiroNaoEditaPost for 6}.
\item[M12\_1-INV-013] TerceiroNaoEditaComent: UNSAT.\newline Escopo: \texttt{check TerceiroNaoEditaComent for 6}.
\item[M12\_1-INV-014] ComentarExigeParticipacaoAtual: UNSAT.\newline Escopo: \texttt{check ComentarExigeParticipacaoAtual for 6}.
\item[M12\_1-INV-015] EditarComentExigeParticipacaoAtual: UNSAT.\newline Escopo: \texttt{check EditarComentExigeParticipacaoAtual for 6}.
\item[M12\_1-INV-016] RemovidoNaoEditaComent: UNSAT.\newline Escopo: \texttt{check RemovidoNaoEditaComent for 6}.
\item[M12\_1-INV-017] RemovidoNaoLe: UNSAT.\newline Escopo: \texttt{check RemovidoNaoLe for 6}.
\item[M12\_1-INV-018] LerExigeParticipacaoOuPapel: UNSAT.\newline Escopo: \texttt{check LerExigeParticipacaoOuPapel for 6}.
\item[M12\_1-INV-019] InativoNaoLe: UNSAT.\newline Escopo: \texttt{check InativoNaoLe for 6}.
\item[M12\_1-INV-020] EdicaoPostCriaHistorico: UNSAT.\newline Escopo: \texttt{check EdicaoPostCriaHistorico for 6}.
\item[M12\_1-INV-021] RemocaoPostCriaHistorico: UNSAT.\newline Escopo: \texttt{check RemocaoPostCriaHistorico for 6}.
\item[M12\_1-INV-022] RemocaoPreservaDocumento: UNSAT.\newline Escopo: \texttt{check RemocaoPreservaDocumento for 6}.
\item[M12\_1-INV-023] HistoricoNuncaRemovido: UNSAT.\newline Escopo: \texttt{check HistoricoNuncaRemovido for 6}.
\item[M12\_1-INV-024] EdicaoComentCriaHistorico: UNSAT.\newline Escopo: \texttt{check EdicaoComentCriaHistorico for 6}.
\item[M12\_1-INV-025] ModeracaoComentCriaHistorico: UNSAT.\newline Escopo: \texttt{check ModeracaoComentCriaHistorico for 6}.
\item[M12\_1-INV-026] EdicaoNaoDesfazModeracao: UNSAT.\newline Escopo: \texttt{check EdicaoNaoDesfazModeracao for 6}.
\item[M12\_1-INV-027] ComentModeradoPreservado: UNSAT.\newline Escopo: \texttt{check ComentModeradoPreservado for 6}.
\item[M12\_1-INV-028] ColegaNaoVeOriginalModerado: UNSAT.\newline Escopo: \texttt{check ColegaNaoVeOriginalModerado for 6}.
\item[M12\_1-INV-029] AutorVeOriginalMarcado: UNSAT.\newline Escopo: \texttt{check AutorVeOriginalMarcado for 6}.
\item[M12\_1-INV-030] AuditorVeOriginal: UNSAT.\newline Escopo: \texttt{check AuditorVeOriginal for 6}.
\item[M12\_1-INV-031] NotificacaoSoDoAlvo: UNSAT.\newline Escopo: \texttt{check NotificacaoSoDoAlvo for 6}.
\item[M12\_1-INV-032] SemAcessoNaoPublicaComRoteiro: UNSAT.\newline Escopo: \texttt{check SemAcessoNaoPublicaComRoteiro for 6}.
\item[M12\_1-INV-033] PublicacaoComRoteiroExigeAcesso: UNSAT.\newline Escopo: \texttt{check PublicacaoComRoteiroExigeAcesso for 6}.
\item[M12\_1-INV-034] PrimeiraExecucaoProduzReceipt: UNSAT.\newline Escopo: \texttt{check PrimeiraExecucaoProduzReceipt for 6}.
\item[M12\_1-INV-035] IdentidadeComandoUnica: UNSAT.\newline Escopo: \texttt{check IdentidadeComandoUnica for 6}.
\item[M12\_1-INV-036] ReusoIncompativelRejeitado: UNSAT.\newline Escopo: \texttt{check ReusoIncompativelRejeitado for 6}.
\item[M12\_1-INV-037] RetryNaoReexecuta: UNSAT.\newline Escopo: \texttt{check RetryNaoReexecuta for 6}.
\item[M12\_1-INV-038] RetryNaoDuplicaFato: UNSAT.\newline Escopo: \texttt{check RetryNaoDuplicaFato for 6}.
\item[M12\_1-INV-039] RevogacaoImpedeCommit: UNSAT.\newline Escopo: \texttt{check RevogacaoImpedeCommit for 6}.
\item[M12\_1-INV-040] TransicoesPreservamCoerencia: UNSAT.\newline Escopo: \texttt{check TransicoesPreservamCoerencia for 6}.
\item[M12\_1-WIT-041] WitnessTurmaAtiva: SAT.\newline Escopo: \texttt{run WitnessTurmaAtiva for 4}.
\item[M12\_1-WIT-042] WitnessCriarPost: SAT.\newline Escopo: \texttt{run WitnessCriarPost for 6}.
\item[M12\_1-WIT-043] WitnessPublicaComRoteiro: SAT.\newline Escopo: \texttt{run WitnessPublicaComRoteiro for 6}.
\item[M12\_1-WIT-044] WitnessEditarPost: SAT.\newline Escopo: \texttt{run WitnessEditarPost for 6}.
\item[M12\_1-WIT-045] WitnessRemoverPost: SAT.\newline Escopo: \texttt{run WitnessRemoverPost for 6}.
\item[M12\_1-WIT-046] WitnessCriarComent: SAT.\newline Escopo: \texttt{run WitnessCriarComent for 6}.
\item[M12\_1-WIT-047] WitnessEditarComent: SAT.\newline Escopo: \texttt{run WitnessEditarComent for 6}.
\item[M12\_1-WIT-048] WitnessModerarComent: SAT.\newline Escopo: \texttt{run WitnessModerarComent for 6}.
\item[M12\_1-WIT-049] WitnessEdicaoNaoDesfazModeracao: SAT.\newline Escopo: \texttt{run WitnessEdicaoNaoDesfazModeracao for 6}.
\item[M12\_1-WIT-050] WitnessLeituraMascarada: SAT.\newline Escopo: \texttt{run WitnessLeituraMascarada for 6}.
\item[M12\_1-WIT-051] WitnessChefeModera: SAT.\newline Escopo: \texttt{run WitnessChefeModera for 6}.
\item[M12\_1-WIT-052] WitnessRemocaoAlunoBloqueia: SAT.\newline Escopo: \texttt{run WitnessRemocaoAlunoBloqueia for 6}.
\item[M12\_1-WIT-053] WitnessRetryAposExecucao: SAT.\newline Escopo: \texttt{run WitnessRetryAposExecucao for 6}.
\item[M12\_1-WIT-054] WitnessClaimAtualSemVinculo: SAT.\newline Escopo: \texttt{run WitnessClaimAtualSemVinculo for 6}.
\item[M12\_1-WIT-055] WitnessReusoIncompativel: SAT.\newline Escopo: \texttt{run WitnessReusoIncompativel for 6}.
\item[M12\_1-WIT-056] WitnessRoteiroAceito: SAT.\newline Escopo: \texttt{run WitnessRoteiroAceito for 6}.
\item[M12\_1-WIT-057] WitnessRetry: SAT.\newline Escopo: \texttt{run WitnessRetry for 4}.
\end{description}
UNSAT para check indica ausência de contraexemplo no escopo declarado; SAT para run indica testemunha encontrada. A evidência não certifica a implementação Firebase/backend.

\par\endgroup

\subsection*{Limites}
Esta evidência prova somente o modelo formal no escopo declarado e não é
homologação da implementação real.
\fussy
